\section{Introducción}

\subsection{Contexto y propósito}

Java no se instala como un programa único. Lo que se instala es una plataforma formada
por tres piezas que conviene distinguir desde el principio, porque confundirlas es el
origen de la mayoría de los problemas durante el proceso:

\begin{itemize}
	\item La \textbf{JVM} (\textit{Java Virtual Machine}) es la máquina virtual que
	ejecuta el \textit{bytecode}. Es lo que hace que un programa compilado en una
	computadora funcione en otra con un sistema operativo distinto.
	\item El \textbf{JRE} (\textit{Java Runtime Environment}) agrupa la JVM más las
	bibliotecas estándar. Con él se pueden \textbf{ejecutar} aplicaciones Java, pero no
	escribirlas.
	\item El \textbf{JDK} (\textit{Java Development Kit}) contiene todo lo anterior más
	las herramientas de desarrollo: el compilador \texttt{javac}, el depurador
	\texttt{jdb}, el empaquetador \texttt{jar}, el generador de documentación
	\texttt{javadoc} y, desde Java~9, el ensamblador de imágenes \texttt{jlink}.
\end{itemize}

La distinción importa porque muchas distribuciones de Linux traen preinstalado
únicamente el JRE. Un equipo así puede abrir cualquier aplicación Java y, sin embargo,
al escribir \texttt{javac Programa.java} responde \texttt{command not found}. Ese es
exactamente el estado inicial de la máquina documentada en este reporte.

\subsection{Por qué hace falta configurar \texttt{JAVA\_HOME}}

Instalar el JDK deja disponibles los comandos \texttt{java} y \texttt{javac} en la
terminal, y para un ejercicio de clase eso podría bastar. El problema aparece en cuanto
se usa cualquier herramienta del ecosistema.

Maven, Gradle, Apache Tomcat, Apache Ant, Elasticsearch, Jenkins y los entornos de
desarrollo como IntelliJ IDEA o Eclipse \textbf{no buscan el compilador en el
\texttt{PATH}}: consultan una variable de entorno llamada \texttt{JAVA\_HOME} que debe
apuntar al \textbf{directorio raíz} de la instalación del JDK. Si esa variable no
existe o apunta al lugar equivocado, esas herramientas fallan con mensajes del estilo
\texttt{JAVA\_HOME is not set and no `java' command could be found}, aunque
\texttt{java -version} funcione perfectamente en esa misma terminal.

Hay una segunda razón, menos evidente pero igual de importante: una computadora puede
tener \textbf{varias versiones del JDK instaladas al mismo tiempo}. Es lo normal en un
equipo de desarrollo, donde un proyecto heredado exige Java~8 y otro requiere Java~21.
En ese escenario, \texttt{JAVA\_HOME} es el mecanismo que declara cuál de todas es la
que debe usarse, y el sistema de alternativas de Debian y Ubuntu
(\texttt{update-alternatives}) es el que resuelve qué binario responde al escribir
\texttt{java}.

\subsection{Objetivos}

De acuerdo con lo solicitado en la práctica:

\begin{enumerate}
	\item \textbf{Identificar} las variables de entorno necesarias para la correcta
	ejecución del Java Development Kit, en particular \texttt{JAVA\_HOME} y su relación
	con \texttt{PATH}.
	\item \textbf{Configurar} dichas variables de forma persistente en un sistema
	GNU/Linux, verificando que la configuración sobreviva al cierre de la terminal y al
	reinicio del equipo.
	\item \textbf{Documentar} el procedimiento completo —verificación del estado
	inicial, desinstalación de versiones previas, instalación, configuración y
	comprobación— con evidencia gráfica de cada etapa.
\end{enumerate}

\subsection{Entorno de trabajo}

\rowcolors{2}{rowblue}{white}
\footnotesize
\renewcommand{\arraystretch}{1.15}
\setlength{\tabcolsep}{4pt}

\begin{longtable}{|>{\raggedright\arraybackslash}p{4.6cm}|>{\raggedright\arraybackslash}p{9.6cm}|}
	\caption{Características del equipo utilizado}
	\label{tab:entorno} \\
	\hline
	\rowcolor{headerblue}
	\headercell{Elemento} & \headercell{Valor} \\
	\hline
	\endfirsthead
	\hline
	\rowcolor{headerblue}
	\headercell{Elemento} & \headercell{Valor} \\
	\hline
	\endhead

	Sistema operativo & Ubuntu 22.04.5 LTS (Jammy Jellyfish) \\ \hline
	Arquitectura & x86\_64 (64 bits) \\ \hline
	Intérprete de comandos & Bash \\ \hline
	Estado inicial de Java & OpenJDK 11 JRE (solo entorno de ejecución, sin compilador) \\ \hline
	Versión a instalar & OpenJDK 21 LTS \\ \hline
	Método de instalación & Gestor de paquetes APT, repositorio \textit{jammy-updates} \\ \hline
	Privilegios requeridos & Cuenta con acceso a \texttt{sudo} \\ \hline
\end{longtable}

\normalsize
\setlength{\tabcolsep}{6pt}

\subsection{Nota sobre la distribución elegida}

La práctica menciona \textit{Oracle Java Development Kit}. En este reporte se documenta
la instalación de \textbf{OpenJDK 21}, y conviene justificar la decisión porque no es
una omisión sino una elección técnica.

Desde Java~7, \textbf{OpenJDK es la implementación de referencia del lenguaje}: el
código fuente del que parte Oracle para construir su propia distribución es el mismo
proyecto OpenJDK. Ambas superan el conjunto de pruebas de compatibilidad TCK, de modo
que un programa compilado con una funciona idénticamente en la otra. Las diferencias
relevantes son de licencia y de soporte comercial, no de comportamiento:

\begin{itemize}
	\item \textbf{OpenJDK} se distribuye bajo licencia GPLv2 con excepción de enlazado,
	es de uso libre en cualquier contexto y llega empaquetado en los repositorios
	oficiales de Ubuntu, lo que simplifica la instalación y las actualizaciones de
	seguridad.
	\item \textbf{Oracle JDK} se distribuye bajo la \textit{Oracle No-Fee Terms and
	Conditions}, que permite el uso gratuito con restricciones, y requiere descargar el
	paquete manualmente aceptando la licencia.
\end{itemize}

Todo lo relativo a \texttt{JAVA\_HOME}, al \texttt{PATH} y al sistema de alternativas
—que es el objetivo de aprendizaje de la práctica— es \textbf{idéntico} en ambos casos;
lo único que cambia es la ruta de instalación. En la sección de conclusiones se retoma
esta comparación.
